\section{操作系统概述}

\subsection{操作系统的概念与功能}

\begin{mydef}{操作系统}{os-def}
操作系统（Operating System, OS）是管理计算机硬件与软件资源的系统软件，也是用户与计算机之间的接口。它提供：
\begin{enumerate}
    \item \textbf{资源管理}：处理器、存储器、I/O设备和文件的分配与调度；
    \item \textbf{抽象}：将底层硬件细节封装为简洁一致的接口（虚拟化）；
    \item \textbf{服务}：为用户和应用程序提供运行环境。
\end{enumerate}
\end{mydef}

\begin{conclusion}{操作系统的四大功能}{os-four-functions}
\begin{tablebox}{OS 四大管理功能}
\centering
\begin{tabular}{|l|l|l|}
\hline
功能 & 管理对象 & 核心问题 \\
\hline
进程管理 & CPU & 进程调度、同步、通信、死锁 \\
内存管理 & 主存 & 分配回收、地址变换、虚拟内存、保护 \\
文件管理 & 外存 & 文件目录、存储空间、共享保护 \\
设备管理 & I/O设备 & 缓冲、分配、驱动、SPOOLing \\
\hline
\end{tabular}
\end{tablebox}
\end{conclusion}

\subsection{操作系统的特征}

\begin{conclusion}{OS 的四个基本特征}{os-features}
\begin{enumerate}
    \item \textbf{并发（Concurrency）}：多个程序在同一时间间隔内同时运行（宏观并行、微观串行）。\textbf{与并行（Parallelism）的区别}：并行是同一时刻真正同时执行（需多核）。
    \item \textbf{共享（Sharing）}：多个进程共享系统资源。分为互斥共享（如打印机）和同时访问（如磁盘）。
    \item \textbf{虚拟（Virtualization）}：将物理实体映射为多个逻辑对应物（虚拟CPU、虚拟内存、虚拟设备）。
    \item \textbf{异步（Asynchrony）}：进程以不可预知的速度推进（走走停停），OS 需保证结果可再现。
\end{enumerate}
\textbf{并发和共享是 OS 最基本的两个特征}，互为存在条件。
\end{conclusion}

\begin{attention}
\textbf{408 常考辨析：并发 vs 并行。} 单核 CPU 通过时间片轮转实现并发，不是并行。多核系统才能实现真正的并行。
\end{attention}

\subsection{操作系统的发展与分类}

\begin{tablebox}{OS 发展历程}
\centering
\begin{tabular}{|l|l|l|}
\hline
阶段 & 特征 & 典型系统 \\
\hline
手工操作 & 人机矛盾突出 & — \\
单道批处理 & 监督程序（Monitor），自动作业过渡 & — \\
多道批处理 & 多道程序并发，资源利用率高，无交互 & IBM OS/360 \\
分时系统 & 时间片轮转，交互性强 & UNIX \\
实时系统 & 严格的响应时间约束 & RTOS \\
\hline
\end{tabular}
\end{tablebox}

\textbf{多道程序设计的核心：} 内存中同时存放多道程序，CPU 在等待 I/O 时切换到另一道程序，提高 CPU 利用率。

\subsection{操作系统内核}

\begin{mydef}{内核}{kernel}
操作系统最核心的部分，常驻内存。分为两大类型：
\begin{itemize}
    \item \textbf{大内核（Monolithic Kernel）}：将 OS 主要功能模块（进程管理、内存管理、文件系统等）都运行在内核态。优点是性能高；缺点是内核庞大，一个模块崩溃可能导致整个系统崩溃。
    \item \textbf{微内核（Microkernel）}：仅将最基本功能（中断处理、进程调度、进程通信）放入内核，其余（文件系统、设备驱动等）作为用户进程运行。优点是可靠性高、易扩展；缺点是频繁的用户态/内核态切换导致性能开销。
\end{itemize}
\end{mydef}

\begin{center}
\begin{tikzpicture}[>=stealth, font=\scriptsize]
% ---- panel titles ----
\node[font=\scriptsize, text=Burgundy, font=\bfseries] at (-3.9,3.9) {大内核（Monolithic Kernel）};
\node[font=\scriptsize, text=Burgundy, font=\bfseries] at (3.9,3.9) {微内核（Microkernel）};
% ---- LEFT: 大内核 ----
\node[draw, rectangle, fill=SoftGray, minimum width=6.6cm, minimum height=0.7cm, rounded corners=2pt] (userL) at (-3.9,2.7) {用户态（User Mode）：用户进程 / 应用};
\node[draw, rectangle, fill=SoftTeal, minimum width=6.6cm, minimum height=3.3cm, rounded corners=3pt] (kernelL) at (-3.9,0.7) {};
\node[font=\scriptsize, text=Navy, align=center] at (-3.9,2.05) {\textbf{内核态（Kernel）}\\ 所有功能装入内核};
\node[draw, rectangle, fill=SoftBlue, minimum width=1.25cm, minimum height=0.55cm] (km1) at (-5.25,1.1) {进程管理};
\node[draw, rectangle, fill=SoftBlue, minimum width=1.25cm, minimum height=0.55cm] (km2) at (-3.9,1.1) {内存管理};
\node[draw, rectangle, fill=SoftBlue, minimum width=1.25cm, minimum height=0.55cm] (km3) at (-2.55,1.1) {文件管理};
\node[draw, rectangle, fill=SoftBlue, minimum width=1.25cm, minimum height=0.55cm] (km4) at (-4.55,0.35) {设备管理};
\node[draw, rectangle, fill=SoftBlue, minimum width=1.25cm, minimum height=0.55cm] (km5) at (-3.25,0.35) {网络};
\node[draw, rectangle, fill=SoftAmber, minimum width=6.6cm, minimum height=0.7cm, rounded corners=2pt] (hwL) at (-3.9,-1.3) {硬件（Hardware）：CPU / 内存 / 设备};
% ---- RIGHT: 微内核 ----
\node[draw, rectangle, fill=SoftGray, minimum width=6.6cm, minimum height=2.4cm, rounded corners=3pt] (userR) at (3.9,1.85) {};
\node[font=\scriptsize, text=Navy, align=center] at (3.9,2.75) {\textbf{用户态（User Mode）}\\ 文件系统 / 设备驱动 / 网络\\ 以用户进程运行};
\node[draw, rectangle, fill=SoftBlue, minimum width=1.3cm, minimum height=0.55cm] (um1) at (2.5,1.35) {文件系统};
\node[draw, rectangle, fill=SoftBlue, minimum width=1.3cm, minimum height=0.55cm] (um2) at (3.9,1.35) {设备驱动};
\node[draw, rectangle, fill=SoftBlue, minimum width=1.3cm, minimum height=0.55cm] (um3) at (5.3,1.35) {网络};
\node[draw, rectangle, fill=SoftTeal, minimum width=6.6cm, minimum height=1.3cm, rounded corners=3pt] (kernelR) at (3.9,0.0) {};
\node[font=\scriptsize, text=Navy, align=center] at (3.9,0.35) {\textbf{内核态（Kernel）}\\ 只保留最基本功能};
\node[draw, rectangle, fill=SoftBlue, minimum width=1.4cm, minimum height=0.5cm] (kk1) at (1.65,-0.25) {进程调度};
\node[draw, rectangle, fill=SoftBlue, minimum width=1.4cm, minimum height=0.5cm] (kk2) at (3.15,-0.25) {中断处理};
\node[draw, rectangle, fill=SoftBlue, minimum width=1.4cm, minimum height=0.5cm] (kk3) at (4.65,-0.25) {进程通信};
\node[draw, rectangle, fill=SoftBlue, minimum width=1.4cm, minimum height=0.5cm] (kk4) at (6.15,-0.25) {内存管理};
\node[draw, rectangle, fill=SoftAmber, minimum width=6.6cm, minimum height=0.7cm, rounded corners=2pt] (hwR) at (3.9,-1.3) {硬件（Hardware）：CPU / 内存 / 设备};
% ---- mode boundaries ----
\draw[dashed, RuleGray] (-7.2,2.35) -- (-0.6,2.35);
\draw[dashed, RuleGray] (0.6,0.65) -- (7.2,0.65);
\node[fill=white, inner sep=1pt, font=\scriptsize, text=MutedText] at (-3.9,2.35) {态切换};
\node[fill=white, inner sep=1pt, font=\scriptsize, text=MutedText] at (3.9,0.65) {态切换};
% ---- comparison note ----
\node[font=\scriptsize, text=MutedText, align=center] at (0,-2.6) {大内核：性能高，但内核庞大，一个模块崩溃可能波及整个系统；\\ 微内核：可靠性高、易扩展，但用户态/内核态切换频繁，性能开销较大};
\end{tikzpicture}
\end{center}

\subsection{中断与异常}

\begin{conclusion}{中断与异常}{interrupt-exception}
中断机制是 OS 实现并发的重要基础，尤其是时钟中断使抢占式调度和分时成为可能。进程切换也可能由系统调用、异常、主动让出 CPU 或进程阻塞触发，因此不能把它绝对化为“没有外部中断就没有进程切换”。

\begin{tablebox}{中断与异常对比}
\centering
\begin{tabularx}{\linewidth}{|>{\centering\arraybackslash}p{2.3cm}|X|X|}
\hline
 & 中断（Interrupt）& 异常（Exception）\\
\hline
来源 & 外设（硬件）& CPU 内部（指令执行）\\
与指令的关系 & 异步（与当前指令无关）& 同步（由当前指令触发）\\
响应方式 & 下一条指令开始前检测 & 指令执行过程中检测\\
示例 & 时钟中断、I/O完成中断 & 缺页、除零、断点\\
分类 & 可屏蔽/不可屏蔽 & 故障(Fault)/陷阱(Trap)/终止(Abort)\\
\hline
\end{tabularx}
\end{tablebox}
\end{conclusion}

\begin{attention}
\textbf{中断 vs 异常（内中断 vs 外中断）的辨析：} 缺页异常是典型的 Fault（故障），处理后通常可重新执行引发异常的指令；系统调用（如 \texttt{int 0x80} 或 \texttt{syscall}）是由程序主动触发的 Trap（陷阱）。除零等同步异常的具体 Fault/Abort 归类依赖体系结构和异常处理约定，不能脱离体系结构一概而论。
\end{attention}

\begin{center}
\begin{tikzpicture}[>=stealth, font=\scriptsize]
  \node[draw,rounded corners=2pt,fill=SoftGray,minimum width=1.7cm,minimum height=0.7cm] (a) at (0,0.4){中断请求};
  \node[draw,rounded corners=2pt,fill=SoftGray,minimum width=1.7cm,minimum height=0.7cm] (b) at (2.3,0.4){中断响应};
  \node[draw,rounded corners=2pt,fill=SoftGray,minimum width=1.7cm,minimum height=0.7cm] (c) at (4.6,0.4){保护现场};
  \node[draw,rounded corners=2pt,fill=SoftGray,minimum width=1.7cm,minimum height=0.7cm] (d) at (6.9,0.4){中断服务};
  \node[draw,rounded corners=2pt,fill=SoftGray,minimum width=1.7cm,minimum height=0.7cm] (e) at (9.2,0.4){恢复现场};
  \draw[->] (a) -- (b); \draw[->] (b) -- (c); \draw[->] (c) -- (d); \draw[->] (d) -- (e);
  \node[font=\scriptsize,black!70] at (9.2,-0.6){返回被中断的程序};
\end{tikzpicture}
\end{center}

\subsection{系统调用}

\begin{mydef}{系统调用}{syscall}
系统调用（System Call）是 OS 提供给用户程序的编程接口，是用户程序主动请求内核服务的规范入口。程序通过体系结构规定的陷入/系统调用机制从用户态进入内核态，执行完毕后返回用户态；中断和其他异常也能使 CPU 进入内核处理流程，但不属于应用主动请求服务的普通接口。

常见系统调用：进程控制（\texttt{fork}, \texttt{exit}, \texttt{wait}）、文件操作（\texttt{open}, \texttt{read}, \texttt{write}）、设备管理（\texttt{ioctl}）、通信（\texttt{pipe}, \texttt{shmget}）。
\end{mydef}

\begin{center}
\begin{tikzpicture}[>=stealth, node distance=0.75cm, font=\scriptsize]
\node[draw, rectangle, fill=SoftGray, minimum width=3.6cm, minimum height=0.6cm] (u) {用户进程（用户态）};
\node[draw, rectangle, fill=SoftAmber, minimum width=3.6cm, minimum height=0.6cm, below of=u] (t) {陷入（trap / syscall）};
\node[draw, rectangle, fill=SoftBlue, minimum width=3.6cm, minimum height=0.6cm, below of=t] (k) {内核：系统调用处理程序（内核态）};
\node[draw, rectangle, fill=SoftGray, minimum width=3.6cm, minimum height=0.6cm, below of=k] (r) {返回用户态，继续执行};
\draw[->] (u) -- (t);
\draw[->] (t) -- (k);
\draw[->] (k) -- (r);
\node[font=\scriptsize, black!60, align=left, right of=t, node distance=3.5cm] {保护现场（保存寄存器、\\ PC、PSW）};
\node[font=\scriptsize, black!60, align=left, right of=k, node distance=3.9cm] {执行请求的服务\\（文件、进程、设备等）};
\node[font=\scriptsize, black!60, align=left, right of=r, node distance=3.7cm] {恢复现场，返回用户态};
\end{tikzpicture}
\end{center}

\subsection{操作系统的体系结构}

\begin{conclusion}{OS 体系结构演化}{os-arch}
\begin{itemize}
    \item \textbf{简单结构}（MS-DOS）：无模块划分，用户程序可直接访问硬件。
    \item \textbf{模块化结构}：按功能划分模块，模块间通过接口调用。
    \item \textbf{层次结构}：底层为硬件相关，高层为用户接口，每层仅调用相邻下层。
    \item \textbf{外核结构（Exokernel）}：内核仅负责资源隔离和保护，资源管理策略交给应用程序。
\end{itemize}
\end{conclusion}

\subsection{408 高频考点速查}

\begin{conclusion}{操作系统概述 — 408常考点}{os1-keypoints}
\begin{enumerate}
    \item \textbf{并发 vs 并行}：单核 CPU 分时 = 并发；多核 = 并行。
    \item \textbf{中断 vs 异常}：来源（外 vs 内）、同步性、典型例子。
    \item \textbf{用户态 vs 内核态}：系统调用/中断/异常触发态切换；特权指令只能在内核态执行。
    \item \textbf{微内核 vs 大内核}：各有哪些 OS 模块在内核中。
    \item \textbf{系统调用与库函数的区别}：系统调用需陷入内核；库函数（如 \texttt{printf}）在用户态执行，内部可能调用系统调用（如 \texttt{write}）。
\end{enumerate}
\end{conclusion}

\newpage
\section{习题}

\subsection{基础过关题}

\begin{enumerate}

\item \textbf{（单项选择题）} 操作系统的主要功能是
\begin{center}
\begin{tabular}{ll}
(A) 实现程序并发执行 & (B) 管理计算机系统中的资源 \\[4pt]
(C) 提供用户界面 & (D) 以上都是
\end{tabular}
\end{center}

\ifshowsolutions\par\noindent\textbf{解析：} 操作系统负责资源管理，为程序并发执行提供机制，并向用户和应用提供接口。\boxed{\textbf{答案：}(D)}\else \answerarea \fi

\item \textbf{（单项选择题）} 以下关于并发和并行的说法，正确的是
\begin{center}
\begin{tabular}{ll}
(A) 单核 CPU 可以并行执行多个进程 \\[4pt]
(B) 并发是指同一时刻多个程序同时运行 \\[4pt]
(C) 同一时刻真正执行多个 CPU 指令流需要多核或多处理器等并行硬件资源 \\[4pt]
(D) 并发和并行是同一概念
\end{tabular}
\end{center}

\ifshowsolutions\par\noindent\textbf{解析：} 并发描述多个任务在同一时间段交替推进；并行描述同一时刻真正同时执行。单核 CPU 可通过切换实现并发，但同一时刻只能执行一个 CPU 指令流。\boxed{\textbf{答案：}(C)}\else \answerarea \fi

\item \textbf{（填空题）} 中断来自 CPU \underline{\qquad} 部，异常来自 CPU \underline{\qquad} 部。系统调用属于 \underline{\qquad}（异常类型）。

\ifshowsolutions\par\noindent\textbf{答案：} 外；内；陷阱（Trap）。外部中断与当前指令异步，异常由指令执行或程序主动触发并与当前指令同步。\else \answerarea \fi

\item \textbf{（简答题）} 简述微内核和大内核的区别，各举一个代表性操作系统。

\ifshowsolutions\par\noindent\textbf{答案：} 大内核把进程管理、内存管理、文件系统和大量驱动等主要服务放在内核态，模块间调用开销低，但故障隔离较弱，代表如传统 UNIX/Linux（Linux 是模块化单体内核）。微内核只保留调度、进程间通信、基本地址空间和中断等最小机制，文件系统和驱动等运行在用户态服务中，隔离和可扩展性较好，但 IPC 与态切换开销更大，代表如 MINIX 3、QNX。具体归类可能因混合内核设计而有边界。\else \answerarea \fi

\item \textbf{（简答题）} 简述操作系统为用户提供的三种用户接口。

\ifshowsolutions\par\noindent\textbf{答案：} 常见教材分类为：命令接口（联机命令/命令解释器和脱机作业控制语言）、图形用户接口，以及供程序请求内核服务的程序接口，即系统调用。前两类面向最终用户，系统调用面向应用程序。\else \answerarea \fi

\end{enumerate}

\subsection{真题风格与综合训练}

\begin{enumerate}[resume]

\item \textbf{（真题风格，单项选择题）} 下列事件中，属于中断的是
\begin{center}
\begin{tabular}{ll}
(A) 用户程序执行系统调用 & (B) 除零错误 \\[4pt]
(C) 缺页异常 & (D) 磁盘 I/O 完成
\end{tabular}
\end{center}

\ifshowsolutions\par\noindent\textbf{解析：} 磁盘 I/O 完成由外设控制器异步发出外部中断；系统调用是陷阱，除零和缺页属于同步异常。\boxed{\textbf{答案：}(D)}\else \answerarea \fi

\item \textbf{（真题风格，单项选择题）} 下列指令中，只能在内核态执行的是
\begin{center}
\begin{tabular}{ll}
(A) 读时钟 & (B) 置时钟 \\[4pt]
(C) 取数指令 & (D) 寄存器清零
\end{tabular}
\end{center}

\ifshowsolutions\par\noindent\textbf{解析：} 修改时钟/定时器会影响系统调度和保护，属于特权操作；普通读取、取数和寄存器运算可在用户态按体系结构权限执行。\boxed{\textbf{答案：}(B)}\else \answerarea \fi

\item \textbf{（真题风格，综合题）} 有 3 个同时到达的作业 A、B、C，各自的 CPU 计算时间和 I/O 时间如下：
\begin{itemize}
    \item A：计算 30ms，I/O 40ms，计算 10ms
    \item B：计算 20ms，I/O 30ms，计算 20ms
    \item C：计算 10ms，I/O 20ms，计算 30ms
\end{itemize}
问：(1) 若三道作业按 A$\to$B$\to$C 顺序单道执行，总时间是多少？(2) 若引入多道程序设计，假定内存可同时容纳三道作业，系统只有一个 CPU，CPU 调度采用非抢占 FCFS，初始顺序为 A、B、C；各作业的 I/O 可彼此并行且可与 CPU 重叠，切换开销忽略不计。画出 CPU 甘特图并计算总时间。

\ifshowsolutions\par\noindent\textbf{答案：}
(1) 单道执行时，CPU 在每个作业 I/O 期间也空闲，总时间为
\[
(30+40+10)+(20+30+20)+(10+20+30)=210\text{ ms}.
\]

(2) CPU 时间线为
\[
\begin{array}{c|ccccccc}
\text{时间段/ms}&0\!\sim\!30&30\!\sim\!50&50\!\sim\!60&60\!\sim\!70&70\!\sim\!80&80\!\sim\!100&100\!\sim\!130\\ \hline
\text{CPU}&A&B&C&\text{空闲}&A&B&C
\end{array}
\]
A、B、C 的 I/O 分别发生在30--70、50--80、60--80 ms。60 ms 时三者都在等待 I/O，CPU 空闲至70 ms；80 ms 时 B、C 同时就绪，按题设初始 FCFS 次序先运行 B。总时间为130 ms。\else \answerarea \fi

\end{enumerate}

\vspace{8pt}
\noindent\textbf{拓展阅读：}
\begin{itemize}
    \item OSTEP 第 1-2 章「操作系统介绍」和「虚拟化」——经典的 OS 全景图和设计哲学；
    \item 陈海波《现代操作系统：原理与实现》第 1 章——现代 OS 的架构演变和 ChCore 实验环境。
\end{itemize}
